% Р11-01 — Смеси и сплавы
% Источник: старая методичка, стр. 90–91

\begin{razbor}{Р11-01}
\Uslovie Имеется два раствора кислоты: 30-процентный и 70-процентный.
Сколько килограммов 30-процентного раствора нужно взять, чтобы, смешав его
со вторым, получить 40 кг 50-процентного раствора?

\Reshenie Пусть взяли $x$ кг первого раствора, тогда второго — $40-x$ кг.
Масса чистой кислоты при смешивании сохраняется (\tref{Ч1.3}):

\begin{pfTablica}{lccc}
Раствор & Масса, кг & Доля кислоты & Кислоты, кг\\
1-й & $x$ & $0{,}3$ & $0{,}3x$\\
2-й & $40-x$ & $0{,}7$ & $0{,}7(40-x)$\\
\pfItog Смесь & $40$ & $0{,}5$ & $20$\\
\end{pfTablica}
\[
  0{,}3x+0{,}7(40-x)=20 \;\Longrightarrow\; -0{,}4x=-8 \;\Longrightarrow\; x=20.
\]
\Otvet{20.}
\end{razbor}
